\clearpage\pagestyle{otMain}\fontsize{11}{15}\selectfont\corrigesChapitre{12}{Pythagore, cercle et droite des milieux}
\begin{blocLivre}{cor-12-1}\begin{corrige}{1}{Reconnaître l’hypoténuse}\demarche{Repérer d’abord le sommet de l’angle droit.}
Le côté opposé à A est $\res{BC}$. L’égalité est $\res{BC^2=AB^2+AC^2}$. Le nom du sommet de l’angle droit détermine le bon côté.
\label{fig-cor-12-1}\figureGeo[corrige]{pyth-01}
\end{corrige}\end{blocLivre}
\begin{blocLivre}{cor-12-2}\begin{corrige}{2}{Un triangle classique}\demarche{Additionner les carrés pour trouver l’hypoténuse.}
Pythagore donne $BC^2=6^2+8^2=100$. Comme BC est positive, $\res{BC=10\text{ cm}}$.
\label{fig-cor-12-2}\figureGeo[corrige]{pyth-02}
\end{corrige}\end{blocLivre}
\begin{blocLivre}{cor-12-3}\begin{corrige}{3}{Le côté manquant}\demarche{Isoler le carré inconnu avant d’extraire la racine.}
$AC^2=13^2-5^2=144$, donc $\res{AC=12\text{ cm}}$. AC est inférieure à l’hypoténuse 13 cm.
\label{fig-cor-12-3}\figureGeo[corrige]{pyth-03}
\end{corrige}\end{blocLivre}
\begin{blocLivre}{cor-12-4}\begin{corrige}{4}{Garder une racine}\demarche{Arrondir uniquement le résultat final.}
$BC^2=16+49=65$, donc $\res{BC=\sqrt{65}\text{ cm}}\approx\res{8{,}1\text{ cm}}$. L’arrondi n’est pas une égalité exacte.
\label{fig-cor-12-4}\figureGeo[corrige]{pyth-04}
\end{corrige}\end{blocLivre}
\begin{blocLivre}{cor-12-5}\begin{corrige}{5}{Le centre du cercle circonscrit}\demarche{Relier l’angle droit au cercle de diamètre l’hypoténuse.}
Comme ABC est rectangle en A, A appartient au cercle de diamètre [BC]. Le centre est donc $\res{O\text{, milieu de }[BC]}$, le rayon vaut $\res{5\text{ cm}}$ et $\res{OA=5\text{ cm}}$ : la médiane issue de l’angle droit mesure la moitié de l’hypoténuse.
\label{fig-cor-12-5}\figureGeo[corrige]{cercle-circonscrit}
\end{corrige}\end{blocLivre}
\begin{blocLivre}{cor-12-6}\begin{corrige}{6}{Tester la réciproque}\demarche{Comparer les deux membres séparément.}
BC est le plus grand côté. $15^2=225$ et $9^2+12^2=225$. Par la réciproque du théorème de Pythagore, ABC est $\res{\text{rectangle en A}}$.
\label{fig-cor-12-6}\figureGeo[corrige]{pyth-05}
\end{corrige}\end{blocLivre}
\begin{blocLivre}{cor-12-7}\begin{corrige}{7}{Utiliser la contraposée}\demarche{Comparer exactement, sans se fier au dessin.}
$BC^2=81$ alors que $AB^2+AC^2=25+49=74$. Ces nombres diffèrent ; d’après la $\res{\text{contraposée du théorème de Pythagore}}$, ABC $\res{\text{n’est pas rectangle}}$.
\label{fig-cor-12-7}\figureGeo[corrige]{pyth-06}
\end{corrige}\end{blocLivre}
\begin{blocLivre}{cor-12-8}\begin{corrige}{8}{La droite des milieux}\demarche{Citer les deux milieux avant d’appliquer le théorème.}
I et J sont les milieux de deux côtés : d’après le théorème de la droite des milieux, $\res{(IJ)\parallel(BC)}$ et $IJ=BC/2=\res{4{,}5\text{ cm}}$.
\label{fig-cor-12-8}\figureGeo[corrige]{milieux-parallele}
\end{corrige}\end{blocLivre}
\begin{blocLivre}{cor-12-9}\begin{corrige}{9}{Un milieu retrouvé}\demarche{Identifier le théorème dont la conclusion est un milieu.}
La droite passe par le milieu de [AB] et elle est parallèle à (BC) : d’après le deuxième théorème des milieux, J est le milieu de [AC]. Donc $\res{AJ=3{,}7\text{ cm}}$.
\label{fig-cor-12-9}\figureGeo[corrige]{milieux-milieu}
\end{corrige}\end{blocLivre}
\begin{blocLivre}{cor-12-10}\begin{corrige}{10}{Un angle à identifier}\demarche{Nommer le sommet opposé au plus long côté.}
Le plus grand côté est AB. $AB^2=100=36+64=AC^2+BC^2$. La réciproque du théorème de Pythagore donne un angle droit $\res{\text{en C}}$, opposé à AB. La position du dessin ne détermine pas l’angle droit.
\label{fig-cor-12-10}\figureGeo[corrige]{pyth-10}
\end{corrige}\end{blocLivre}
\begin{blocLivre}{cor-12-11}\begin{corrige}{11}{Une échelle contre un mur}\demarche{Justifier l’angle droit par le mur et le sol.}
ABC est rectangle en A, car le mur est perpendiculaire au sol. $AC^2=5^2-1{,}4^2=23{,}04$, donc $\res{AC=4{,}8\text{ m}}$. La hauteur reste inférieure à la longueur de l’échelle.
\label{fig-cor-12-11}\figureGeo[corrige]{pyth-07}
\end{corrige}\end{blocLivre}
\begin{blocLivre}{cor-12-12}\begin{corrige}{12}{Un triangle dans un demi-cercle}\demarche{Utiliser le cercle pour obtenir l’angle droit, puis Pythagore.}
A appartient au cercle de diamètre [BC] et diffère de B et C : ABC est $\res{\text{rectangle en A}}$. Le diamètre vaut $BC=13$ cm. Pythagore donne $AC^2=13^2-5^2=144$, donc $\res{AC=12\text{ cm}}$.
\label{fig-cor-12-12}\figureGeo[corrige]{demi-cercle}
\end{corrige}\end{blocLivre}
\begin{blocLivre}{cor-12-13}\begin{corrige}{13}{Un rectangle sans équerre}\demarche{Utiliser une propriété caractéristique du rectangle par ses diagonales.}
Les diagonales [AC] et [BD] ont le même milieu O : ABCD est un parallélogramme. Elles ont la même longueur, celle d’un diamètre : un parallélogramme dont les diagonales ont même longueur est $\res{\text{un rectangle}}$. On n’a utilisé que le compas et la règle : $\res{\text{aucun angle droit n’est tracé}}$.
\label{fig-cor-12-13}\figureGeo[corrige]{rectangle-compas}
\end{corrige}\end{blocLivre}
\begin{blocLivre}{cor-12-14}\begin{corrige}{14}{Une hauteur de triangle}\demarche{Justifier la perpendicularité avant Pythagore.}
Dans un triangle isocèle, la médiane [CH] est aussi une hauteur : ACH est rectangle en H. AH = 5 cm et $CH^2=13^2-5^2=144$, donc CH = 12 cm. L’aire vaut $10\times12/2=\res{60\text{ cm}^2}$.
\label{fig-cor-12-14}\figureGeo[corrige]{pyth-iso}
\end{corrige}\end{blocLivre}
\begin{blocLivre}{cor-12-15}\begin{corrige}{15}{Le théorème de Varignon}\demarche{Utiliser la diagonale [AC] comme intermédiaire dans deux triangles.}
Dans le triangle ABC, I et J sont des milieux : $(IJ)\parallel(AC)$ et $IJ=AC/2$. Dans le triangle ACD, L et K sont des milieux : $(LK)\parallel(AC)$ et $LK=AC/2$. Les côtés opposés [IJ] et [LK] sont donc parallèles et de même longueur : IJKL est $\res{\text{un parallélogramme}}$, quelle que soit la forme de ABCD.
\label{fig-cor-12-15}\figureGeo[corrige]{varignon}
\end{corrige}\end{blocLivre}
